\section{Mag (Magmas Finitos)}
A categoria de magmas finitos é a estrutura algébrica mais geral que se pode considerar: impõe apenas a existência de uma operação binária, sem qualquer axioma adicional como associatividade ou comutatividade. Por essa razão, serve de ponto de partida para categorias mais ricas como monoides ou grupos.

\subsection{Objetos}
Os objetos de $\mathbf{Mag}$ são magmas finitos, ou seja, conjuntos finitos $M$ equipados com uma operação binária $\cdot\,: M \times M \to M$ sem axiomas adicionais. Em MATLAB, um magma é representado por uma \texttt{struct} que contém o conjunto de elementos e a operação binária definida como uma função anónima.

\begin{lstlisting}[caption={Representação de um magma finito}]
% Definicao de um magma: conjunto {1,2,3} com operacao mod 3
Mag = struct();
Mag.set = {1, 2, 3};           % Conjunto finito
Mag.op  = @(a, b) mod(a + b, 3); % Operacao binaria (sem axiomas adicionais)
\end{lstlisting}

\subsection{Morfismos}
Os morfismos em $\mathbf{Mag}$ são aplicações que preservam a estrutura do magma. Dados dois magmas $(M, \cdot)$ e $(N, \ast)$, uma aplicação $f: M \to N$ é um morfismo se e só se
\[
    f(a \cdot b) = f(a) \ast f(b), \quad \forall\, a, b \in M.
\]
Não são impostos axiomas adicionais na operação binária. A verificação computacional percorre todos os pares de elementos do domínio e compara as imagens pela operação com as operações nas imagens.

\begin{lstlisting}[caption={Verificação de morfismo entre magmas}]
% Construtor de verificador de morfismo entre dois magmas M1 e M2
MorphismMag = @(M1, M2) @(f) ...
    all(arrayfun(@(a, b) M2.op(f(a), f(b)) == f(M1.op(a, b)), ...
        M1.set, M1.set));

% Funcao auxiliar equivalente (forma explícita)
function ok = check_morphism(M1, M2, f)
    ok = all(arrayfun(@(a, b) M2.op(f(a), f(b)) == f(M1.op(a, b)), ...
             M1.set, M1.set));
end
\end{lstlisting}
